%!TEX encoding = UTF-8 Unicode
% LaTeX Template for CV/Resume
%-------------------------------------------------------------------------------
% LuaLaTeX compilation needed
%-------------------------------------------------------------------------------
% Huanyu Zhou <spica.h.zhou@gmail.com>
% https://spica-vir.github.io/about/
%-------------------------------------------------------------------------------

% configure

\documentclass[a4paper,11pt]{article}
\usepackage{academicCV_en}

% contacts

\myname{A Good Man}
\myposition{Scientific Researcher}
\myaddress{Distinguished Department, A Famous University, City 1 000001, State A}

\mytel{0123456789}
\myemail{goodman@email.com}
\mysite{this.url.does.not.exist/} % Personal website, without 'https://', optional

% Some self-defined commands to make CV looks better

% Highlight my name in author list
\newcommand{\MeAsAuthor}{\textbf{\underline{Goodman A}}}
\newcommand{\MeAsCorrespondingAuthor}{\textbf{\underline{Goodman A*}}}

% Add comments (editor's pick, under review, etc.) to a paper entry in the publication list
\newcommand{\PaperComment}[1]{\textbf{\underline{(#1)}}}

\begin{document}

%-------------------------------------------------------------------------------
% Contact Info
%-------------------------------------------------------------------------------

\makeheader % basic information defined above

% academic social media
%
% The following commands are defined in a similar way to '\href'.
% For other social media, commands need to be defined in 'my_CVstyle.sty'
% \myGoogleScholar
% \myORCID
% \myResearchGate
% \myAcademia
% \myGitHub
% \myLinkedIn
% \myTwitter
%
\begin{acmedia}
	\myGoogleScholar{https://scholar.google.com/}{Google Scholar}
	\EntrySeparator
	\myORCID{https://orcid.org/}{ORCID}
	\EntrySeparator
	\myGitHub{https://github.com/}{GitHub}
	\EntrySeparator
	\myLinkedIn{https://www.linkedin.com/}{LinkedIn}
\end{acmedia}

%-------------------------------------------------------------------------------
% Employment History
%-------------------------------------------------------------------------------

\section{Employment History}

% an entry of education / employment
% Usage: \experience{<position>}{<date>}{<organization>}{<location>}{<description(optional)>}

\experience
	{Scientific Researcher}
	{2025 -- Now}
	{Distinguished Department, A Famous University}
	{City 1, State A}
	{\begin{itemize}
		\item[-] Whatever explains this position
	\end{itemize}}

%-------------------------------------------------------------------------------
% Other Positions and Membership of Professional Bodies
%-------------------------------------------------------------------------------

\section{Other Positions and Membership of Professional Bodies}

\experience
	{Membership / Fellowship}
	{2021 -- Now}
	{A society}
	{City 2, State B}{}

\experience
	{A part-time position}
	{2022 -- 2023}
	{Advisory board of a company}
	{City 3, State C}{}

%-------------------------------------------------------------------------------
% Education
%-------------------------------------------------------------------------------

\section{Education}

\experience
	{Doctor of Philosophy}
	{2021 -- 2025}
	{Distinguished Department, A Famous University}
	{City 1, State A}
	{\begin{itemize}
		\item[-] Supervisor: Prof. Brilliant
	\end{itemize}}

\experience
	{Master of Science}
	{2020 -- 2021}
	{Celebrated Department, Yet Another Famous University}
	{City 2, State B}
	{\begin{itemize}
		\item[-] Supervisor: Prof. Genius
		\item[-] GPA: \textbf{100}/100, Rank: \textbf{1}/40
	\end{itemize}}

\experience
	{Bachelor of Science}
	{2016 -- 2020}
	{Renowned College, An Equally Famous University}
	{City 3, State C}
	{\begin{itemize}
		\item[-] FYP supervisor: Prof. Talented
		\item[-] GPA: \textbf{100}/100, Rank: \textbf{top 1\%}
	\end{itemize}}

%-------------------------------------------------------------------------------
% Research interests
%-------------------------------------------------------------------------------
\section{Research Interests}

\begin{text}
\begin{itemize}[leftmargin=1em]
	\item Research topic A that makes you famous;
	\item Research topic B that makes you rich;
	\item Research topic C that makes you smart.
\end{itemize}
\end{text}

%-------------------------------------------------------------------------------
% Awards and Fellowships
%-------------------------------------------------------------------------------

\section{Awards and Fellowships}

% an entry of fellowship / honor / award etc. 
% Usage: \award{<name>}{<date>}{<event>}{<description(optional)>}

\award
	{A Big Fellowship}
	{2026 -- Now}
	{A research organization}
	{Some details about this fellowship}

\award
	{An Award}
	{2024}
	{International Conference}{}

%-------------------------------------------------------------------------------
% Grants
%-------------------------------------------------------------------------------

\section{Research Grants}

\subsection{Role: Principal Investigator}

% an entry of research grants
% Usage: \grant{<name>}{<date>}{<source>}{<amount(optional)>}{<description(optional)>}

\grant
	{The fundamental research on an important topic}
	{2025 -- 2030}
	{Funding Body A (code: AB123456/X)}
	{£ 1,000,000}{}

\subsection{Role: Key Personnel}

\grant
	{The application of a new theory}
	{2021 -- 2024}
	{A Collaborative Company (code: CD987654/Y)}
	{£ 500,000}{
		\begin{itemize}
			\item[-] Collaborators: Dr Smart (Institute A), Prof. Intelligent (University B)
		\end{itemize}
	}

\subsection{Compute and Travel Grants}

\grant
	{National Supercomputer}
	{2022 -- Now}
	{The HPC Center, project code 01234)}
	{1 M CPUh}{}

\grant
	{Conference Travel Grant}
	{2023}
	{Conference Organizer}
	{£500}{}

%-------------------------------------------------------------------------------
% List of Publications
%-------------------------------------------------------------------------------

\section{Publications}

% Publication list environment
% Usage: \begin{paperlist}[ascending=<true/false>]\end{paperlist}

\begin{paperlist}[ascending=true]
	% Preprint entry, must be used within paperlist environment
	% Usage: \preprint{<author>}{<title>}{<journal/server>}{<status>}{<year>}{DOI(optional, without https://doi.org/)}{<other comments(optional)>}
	\preprint
		{\MeAsCorrespondingAuthor, Collaborator A, Collaborator B}
		{Preprint for a ground-breaking research on topic A}
		{arXiv}
		{Preprint}
		{2026}{}{}

	% Journal paper entry, must be used within paperlist environment
	% Usage: \journal{<author>}{<title>}{<journal>}{<year>}{<volume(optional)>}{<issue(optional)>}{<page(optional)>}{DOI(optional, without https://doi.org/)}{<other comments(optional)>}
	\journal
		{Collaborator A, \MeAsAuthor, Collaborator B}
		{Journal paper for a theoretical research on topic B}
		{Peer-reviewed Journal}
		{2026}
		{00}
		{00}
		{0001--0009}{}{\PaperComment{Cover Article}}

	% Conference paper entry, must be used within paperlist environment
	% Usage: \conference{<author>}{<title>}{<conference>}{<location>}{<date>}{<URL(optional)>}{<other comments(optional)>}
	\conference
		{\MeAsCorrespondingAuthor, Collaborator A}
		{Conference paper for an experimental research on topic C}
		{International Conference on Experimental Research}
		{City, State}
		{11th Dec. 2025}{}{}

	% Book chapter entry, must be used within paperlist environment
	% Usage: \chapter{<author>}{<chaptertitle>}{<editor(optional)>}{<booktitle>}{<edition(optional)>}{<location(optional)>}{<publisher>}{<year>}{<page>}{<URL(optional)>}{<other comments(optional)>}
	\chapter
		{\MeAsAuthor}
		{Fundamental Theory on Topic D}
		{Mr Editor, Ms Editor}
		{Introduction to Topic D}
		{1st}
		{Location}
		{Publisher}
		{2020}
		{16-32}{}{}

\end{paperlist}

%-------------------------------------------------------------------------------
% List of Presentations
%-------------------------------------------------------------------------------

\section{Presentations}

\subsection{Invited Presentations}
% Presentation list environment
% Usage: \begin{prelist}\end{prelist}
\begin{prelist}
	% keynote presentation entry, must be used within prelist environment
	% Usage: \keynote{<title>}{<conference>}{<location>}{<date>}{<URL(optional)>}{<other comments(optional)>}
	\keynote
		{The progress review of research topic A}
		{International Conference on Topic A}
		{City, State}
		{2024}{https://www.youtube.com/}{}

	% guest presentation entry, must be used within prelist environment
	% Usage: \guest{<title>}{<conference>}{<location>}{<date>}{<URL(optional)>}{<other comments(optional)>}
	\guest
		{My research on topic B}
		{Departmental Seminar}
		{Department, University}
		{2023}{}{}
\end{prelist}

\subsection{Contributed Presentations}
\begin{prelist}
	% oral presentation entry, must be used within prelist environment
	% Usage: \oral{<title>}{<conference>}{<location>}{<date>}{<URL(optional)>}{<other comments(optional)>}
	\oral
		{The presentation of some preliminary results}
		{Regional Conference on Topic C}
		{City, State}
		{2021}{}{}

	% poster presentation entry, must be used within prelist environment
	% Usage: \poster{<title>}{<conference>}{<location>}{<date>}{<URL(optional)>}{<other comments(optional)>}
	\poster
		{To exchange ideas on topic D}
		{Discussion meeting by a society}
		{Academic Society}
		{2022}{}{}

\end{prelist}

%-------------------------------------------------------------------------------
% Teaching and Mentoring experiences
%-------------------------------------------------------------------------------
\section{Teaching and Mentoring Experiences}

\subsection{Lectures}

\experience
	{Introduction to Research}
	{2025 -- Now}
	{Lecturer, A Famous University}
	{}{}

\experience
	{Experimental Lab}
	{2021 -- 2025}
	{Teaching Assistant, Yet Another Famous University}
	{}{}

\subsection{Research Projects}

\experience
	{An interesting research on topic A}
	{2026 -- Now}
	{PhD, A Famous University}
	{Mr Student}
	{}

\experience
	{An even more interesting research on topic B}
	{2024 -- 2026}
	{MSi, A Famous University}
	{Ms Student}
	{}

\end{document}
